\chapter{Reglas de negocio}\label{cap:reglas}
Este capítulo concentra las reglas que el sistema aplica de forma automática. Cada una se implementa en un módulo del código (\cref{tab:modulos}) para que pueda revisarse y modificarse en un solo lugar.

\section{Estados de los reportes}
Los reportes diarios y mensuales siguen el mismo ciclo (\cref{fig:estados-reporte}). Solo el Residente los valida y decide una sola vez; el estatus del reporte se deriva siempre de su decisión.

\begin{figure}[h]
\centering
\begin{tikzpicture}[node distance=1.2cm and 1.5cm,
  est/.style={rectangle,rounded corners=8pt,draw=#1,fill=#1!8,text=black,font=\sffamily\small\bfseries,minimum width=2.3cm,minimum height=0.8cm},
  flecha/.style={-{Stealth[length=2.4mm]},line width=0.6pt,draw=black!60},
  et/.style={font=\sffamily\scriptsize,text=black!70,fill=white,inner sep=1.5pt}]
  \node[est=gristexto] (b) {Borrador};
  \node[est=azulinfo,right=of b] (p) {Pendiente};
  \node[est=verdeok,right=2.4cm of p] (a) {Aprobado};
  \node[est=rojo,below=1.5cm of a] (x) {Rechazado};
  \draw[flecha] (b) -- node[et,above]{enviar} (p);
  \draw[flecha] (p) -- node[et,above]{aprueba} (a);
  \draw[flecha] (p) -- node[et,sloped,below]{rechaza} (x);
  \draw[flecha] (x.west) -| node[et,pos=0.2,below]{corregir y reenviar} ([xshift=-4mm]p.south);
\end{tikzpicture}
\caption{Estados de un reporte diario o mensual}\label{fig:estados-reporte}
\end{figure}

\begin{itemize}
  \item Un borrador solo lo ven su autor y los administradores de la obra.
  \item Si el reporte ya tenía decisiones y su autor lo corrige, vuelve a \estatus{Pendiente} y las decisiones se reinician.
  \item Un reporte aprobado ya no se edita. Al aprobarse un reporte mensual, su consolidado queda congelado.
  \item El Administrador puede cambiar cualquier decisión, siempre con motivo (CU-14, CU-15).
\end{itemize}

\section{Estados de las estimaciones}
Las estimaciones se validan en dos etapas sucesivas (\cref{fig:estados-estimacion}): primero el Supervisor y después el Residente.

\begin{figure}[h]
\centering
\begin{tikzpicture}[node distance=1.2cm and 1.25cm,
  est/.style={rectangle,rounded corners=8pt,draw=#1,fill=#1!8,text=black,font=\sffamily\small\bfseries,minimum width=2.2cm,minimum height=0.8cm,align=center},
  flecha/.style={-{Stealth[length=2.4mm]},line width=0.6pt,draw=black!60},
  et/.style={font=\sffamily\scriptsize,text=black!70,fill=white,inner sep=1.5pt,align=center}]
  \node[est=gristexto] (b) {Borrador};
  \node[est=azulinfo,right=of b] (p) {Pendiente};
  \node[est=ambar,right=of p] (r) {En revisión};
  \node[est=verdeok,right=of r] (a) {Aprobado};
  \node[est=guinda,below=1.6cm of p] (o) {Con\\observaciones};
  \node[est=rojo,below=1.6cm of r] (x) {Rechazado};
  \draw[flecha] (b) -- node[et,above]{enviar} (p);
  \draw[flecha] (p) -- node[et,above]{Supervisor\\aprueba} (r);
  \draw[flecha] (r) -- node[et,above]{Residente\\aprueba} (a);
  \draw[flecha] (p) -- node[et,left]{Supervisor} (o);
  \draw[flecha] (p) -- node[et,pos=0.35,right]{Supervisor\\rechaza} (x);
  \draw[flecha] (r) -- node[et,right]{Residente\\rechaza} (x);
  \draw[line width=0.6pt,draw=black!60] (o.south) -- ++(0,-0.55) coordinate (o1);
  \draw[line width=0.6pt,draw=black!60] (x.south) -- ++(0,-0.55) -- (o1);
  \draw[flecha] (o1) -| node[et,pos=0.25,below]{corregir y reenviar} (b.south);
\end{tikzpicture}
\caption{Estados de una estimación}\label{fig:estados-estimacion}
\end{figure}

\begin{itemize}
  \item Solo puede haber una estimación abierta (no aprobada) por contrato, para que los acumulados y las retenciones cuadren.
  \item Mientras está \estatus{Pendiente} o \estatus{En revisión}, sus volúmenes y documentos no se modifican.
  \item \estatus{Con observaciones} regresa la estimación al Superintendente sin que cuente como rechazo; ambos casos exigen comentario.
  \item Al reenviarla, las decisiones previas se reinician y queda constancia en el historial.
\end{itemize}

\section{Cálculos de la estimación}
Sean $I$ el importe de los trabajos del periodo, $a$, $t$ y $r$ los porcentajes de anticipo, IVA y retención por derechos del contrato, $R_{\text{atr}}$ la retención por atraso y $D$ la devolución de retenciones anteriores. Entonces:
\begin{align*}
I &= \textstyle\sum_{i} q_i \, p_i & &\text{(cantidad ejecutada por precio unitario de cada concepto)}\\
A &= I \cdot a & &\text{amortización del anticipo (\reg, art.~143)}\\
\mathit{IVA} &= (I - A)\cdot t\\
R_{\text{der}} &= I \cdot r & &\text{retención por derechos (\reg, art.~128)}\\
\text{Neto} &= I - A + \mathit{IVA} - R_{\text{der}} - R_{\text{atr}} + D
\end{align*}
Los totales los calcula la base de datos con disparadores, de modo que no pueden quedar inconsistentes con los renglones.

\subsection{Retención económica por atraso}
Con el avance físico programado $P$ y el real $F$ a la fecha de corte (último reporte diario aprobado), el monto del contrato $M$ y el porcentaje de retención por atraso $\rho$ (5\,\% por omisión):
\[ R_{\text{req}} = \max(0,\; M \cdot (P - F)) \cdot \rho \]
La retención de la estimación es lo que falta para llegar a $R_{\text{req}}$ respecto de lo ya retenido; si la obra se regulariza, la diferencia se devuelve en la siguiente estimación (\ley, art.~46 Bis). No se aplica cuando el Superintendente justifica un atraso no imputable al contratista.

\begin{aclaracion}[Ejemplo]
Contrato de \$1{,}000{,}000.00 con anticipo de 30\,\%, IVA de 16\,\% y retención por derechos de 0.5\,\%. En enero se estiman 10\,m\textsuperscript{3} a \$1{,}000.00: $I = \$10{,}000.00$, $A = \$3{,}000.00$, $\mathit{IVA} = \$1{,}120.00$ y $R_{\text{der}} = \$50.00$. Al corte el programa marca 30\,\% y el avance real es 20\,\%: $R_{\text{atr}} = \$1{,}000{,}000.00 \times \text{10\,\%} \times \text{5\,\%} = \$5{,}000.00$. Neto a pagar: \$3{,}070.00. Si en febrero la obra alcanza lo programado, la siguiente estimación devuelve los \$5{,}000.00.
\end{aclaracion}

\section{Avance físico, desviación y atraso en días}
\begin{itemize}
  \item El programa de obra registra el avance programado acumulado en fechas clave; entre dos fechas se interpola en línea recta y, antes de la primera, el avance parte de 0\,\% en la fecha de inicio oficial.
  \item El avance real es el del último reporte diario aprobado.
  \item La desviación se expresa en puntos porcentuales (real menos programado) y en días: los días entre la fecha del reporte y la fecha en que el programa alcanzaba el avance real.
  \item Con un día o más de atraso, el reporte diario exige documentar la causa y si se solicita reprogramación o retención (flujo 4a de CU-03).
  \item El avance financiero es la suma de las estimaciones aprobadas entre el monto del contrato.
\end{itemize}

\section{Plazos de las estimaciones}
\begin{itemize}
  \item El periodo de una estimación no puede ser mayor a un mes ni traslaparse con el anterior (\ley, art.~54).
  \item El sistema muestra si se presentó dentro de los 6 días naturales siguientes al corte y cuántos días quedan de los 15 de revisión; al vencer avisa, pero no la autoriza de forma automática.
\end{itemize}

\section{Avisos por correo}
Los avisos se envían por la API de Brevo con una versión en HTML y otra en texto. Un fallo en el correo nunca interrumpe la operación del usuario.

\begin{table}[h]
\centering\small
\caption{Avisos por correo}\label{tab:avisos}
\rowcolors{2}{grisfondo}{white}
\begin{tabularx}{\linewidth}{@{}>{\raggedright\arraybackslash}p{6cm}X@{}}
\toprule
\sffamily\bfseries Evento & \sffamily\bfseries Se avisa a\\\midrule
Reporte diario o mensual enviado a revisión & Residente\\
Reporte aprobado o rechazado & Supervisor (y Directivos, si se aprueba un mensual)\\
Nota de bitácora registrada o corregida & Firmantes que faltan\\
Nota firmada o rechazada & Autor de la nota\\
Estimación enviada & Supervisor\\
Estimación aprobada por el Supervisor & Superintendente y Residente\\
Estimación con observaciones o rechazada por el Supervisor & Superintendente\\
Decisión del Residente sobre una estimación & Superintendente y Supervisor\\
Cambio de estatus por el Administrador & Autor y validadores involucrados\\
\bottomrule
\end{tabularx}
\end{table}
